% ATS_Repair_Rigidity_v0.3_en.tex — Gap Series Part 6
% Corrective R18 manuscript, 2 October 2026. Source-bound claim limits.
\input{glyphtounicode}
\pdfgentounicode=1
\documentclass[11pt,a4paper]{article}
\usepackage{cmap}
\usepackage[utf8]{inputenc}
\usepackage[T1]{fontenc}
\usepackage[american]{babel}
\usepackage{lmodern}
\usepackage{microtype}
\emergencystretch=3em
\usepackage{amsmath,amssymb,amsthm}
\usepackage{booktabs}
\usepackage{tabularx}
\usepackage[colorlinks=true,linkcolor=blue,citecolor=blue,urlcolor=blue]{hyperref}
\usepackage{xurl}

\renewcommand*{\HyperDestNameFilter}[1]{en.#1}

\hypersetup{
  pdftitle={What the Product Formula Allows: A Rigidity Theorem and an Effective-Weight Test for the ATS Height Subchain},
  pdfauthor={Lukas Geiger},
  pdfsubject={Number Theory, Arithmetic Geometry, Teichmüller Theory, abc Conjecture},
  pdfkeywords={ATS, Arithmetic Teichmüller Space, Product Formula, Rigidity, Projective Height, Effective Weights, Joshi},
  pdfcreator={LaTeX with hyperref},
  pdfproducer={pdfTeX},
  bookmarksnumbered=true,
  pdfpagemode=UseOutlines,
  pdfstartview={FitH},
  unicode=true
}

\newtheorem{theorem}{Theorem}[section]
\newtheorem{lemma}[theorem]{Lemma}
\newtheorem{proposition}[theorem]{Proposition}
\newtheorem{corollary}[theorem]{Corollary}
\theoremstyle{definition}
\newtheorem{definition}[theorem]{Definition}
\theoremstyle{remark}
\newtheorem{remark}[theorem]{Remark}

\newcommand{\OL}{\mathcal{O}_L}
\newcommand{\Geff}{\mathcal{G}_{\mathrm{eff}}}

\title{What the Product Formula Allows:\\A Rigidity Theorem and an Effective-Weight Test\\for the ATS Height Subchain}
\author{Lukas Geiger\\\small Independent Researcher, Bernau, Germany\\\small ORCID: 0009-0005-7296-1534}
\date{2 October 2026 --- Version 0.3 (draft)}

\begin{document}
\maketitle

\begin{abstract}
Part~5 of this series \cite{part5} reports a source-bound ATS substance audit;
this corrective Part~6 states the exact scope of its effective-weight test.
For a fixed number field $L$, its fixed standard absolute values, and the
full product formula tested on every $x\in L^\times$, a family of real
exponents $(\gamma_v)$ must be constant across all places, independently of
the class number. Finite-fiber redistribution gives the same conclusion for
effective weights. A prescribed $j^2$-weight on a nonempty place set admits
the constant global extension; only unequal final place weights are excluded.
For global quantity families with a proved linear effective-weight
representation, the product formula leaves at most a $y$-dependent global
scalar. Mere dependence on the weights is insufficient for class membership.
Under ATS~II($\tfrac12$) (5.3.3), direct substitution identifies the weighted
degree on $L^\times$ and the projective height on $\mathbb P^n(L)$ with their
classical values; the source's stabilization on $L^\times$ is then constant
term by term. The extension to $\mathbb P^n(\overline L)$, the ATS~III
line-bundle height, and the ATS~IV volume apparatus require separate bridges.
Log-volume and measure membership or nonmembership in the linear class is
unclassified here. Coordinate hyperplanes and their pullbacks are distinct
objects. No global closure of the ATS repair program, proof or refutation of
abc/IUT, or mathematical certification of the whole ATS series is claimed.
\end{abstract}

\clearpage
\tableofcontents
\clearpage

%=====================================================================
\section{Mandate, Scope, and Source Versions}\label{sec:mandate}

This paper is Part~6 of the abc Gap Series. Part~1 \cite{part1} localized the
quantitative gap in the transport step of Corollary~3.12 of IUTchIII
\cite{mochizukiIUT}; Part~2 \cite{part2} supplied a form criterion;
Part~3 \cite{part3} distilled a checkable admissibility
standard; Part~4 \cite{part4} audited the ATS series
\cite{joshiI,joshiII,joshiIIhalf,joshiIII,joshiIV} at ledger level (v1-F). The
present paper reports the mathematical core of the repair program: a
rigidity theorem, a factorization instrument, and their source-bound
application to the \emph{height subchain} of the ATS texts.

\emph{Scope discipline.} ``The audited passages'' means: ATS~II($\tfrac12$)
\S\S5.3, 5.10, 8.2--8.7 and ATS~IV \S1.5, in the exact arXiv versions cited
in the bibliography (II($\tfrac12$): v12; IV: v2; both revised 24 February
2025); ATS~III v4, Remark~8.9.2(2), is a separately typed
application target. These are publicly posted, versioned preprints, not peer-reviewed
publications; ATS~IV carries the header ``Preliminary version for comments.''
This paper does not audit the full ATS series, does not classify its volume
and measure apparatus (see Case~C of Section~\ref{sec:consequences}), and
makes no claim about structures outside the audited passages. Ground rules as
in Parts~4--5: no authority argument, no truth verdict, verbatim anchors;
findings about later program work appear only as dated external status
notices (Section~\ref{sec:notices}), separate from the proofs of this paper.

%=====================================================================
\section{Types and Operational Glossary}\label{sec:types}

Fix a number field $L$ and its full set $V_L$ of archimedean and finite
places. The local inputs and global quantity families have different types:
\begin{itemize}
\item $b_v$ is the fixed standard absolute value: $|\sigma(x)|$ at a real
place, $|\sigma(x)|^2$ at a complex place, and
$N\mathfrak p^{-v_{\mathfrak p}(x)}$ at a finite place. Thus
$\prod_v b_v(x)=1$ for every $x\in L^\times$.
\item $\rho_{y,v}=b_v^{\,a_{y,v}}$, with $a_{y,v}>0$, is the equivalent
source absolute value on $L_v$. We abbreviate the source datum by $y$.
\item $\lambda_{y,v}$ is a normalization exponent. In our notation,
ATS~II($\tfrac12$) (5.3.3) \cite{joshiIIhalf} says
$b_v=\rho_{y,v}^{\,\lambda_{y,v}}$. This is a paraphrase with renamed
symbols, not a verbatim quotation of the source's $\alpha/\xi$ notation.
\item $\beta_{y,v}=a_{y,v}\lambda_{y,v}$ is the effective weight.
The pointwise identity (5.3.3) implies $\beta_{y,v}=1$ for every $v$;
it also implies $\lambda_{y,v}>0$.
\end{itemize}
The rigidity theorem constrains effective weights, not the two factors
separately. Coordinate variation therefore need not change metric values.

For a precise coordinate comparison use the common finite-support space
$\mathbb R^{(V_L)}$ and
\[
H_y=\{t\in\mathbb R^{(V_L)}:\sum_v\lambda_{y,v}t_v=0\},
\qquad \ell_y(x)=(\log\rho_{y,v}(x))_v.
\]
The vector $\ell_y(x)$ has finite support for $x\in L^\times$.
The normalization identity gives $\ell_y^{-1}(H_y)=L^\times$.
Nonproportional coefficient families define different kernels in this
common space, while these pullbacks coincide. This conditional linear
algebra observation neither proves nonproportionality for a particular
source parameter pair nor verifies the entirety of Theorem~5.10.1.

A family \emph{factors through effective weights} here only when it has the
linear representation in Definition~\ref{def:geff}.
A \emph{classical} value is one obtained from the fixed standard family
on the stated arithmetic domain. A \emph{place-selective} weight family
has unequal final weights, rather than a single global scalar.

\begin{table}[htbp]
\centering\small
\caption{Local inputs, coordinates, and global quantity families.}\label{tab:type_glossary}
\begin{tabularx}{\textwidth}{@{}lXX@{}}
\toprule
Object & Type and dependence & Relation to $\Geff$\\
\midrule
$b_v$ & Fixed local reference & Input, not a global family\\
$\rho_{y,v}$ & Possibly varying local absolute value & Input\\
$\lambda_{y,v}$ & Normalization coordinate & Parameter\\
$\beta_{y,v}$ & Effective coefficient; identically $1$ under (5.3.3) & Coefficient\\
$H_y$ & Coordinate kernel; variation conditional on coefficients & Coordinate object\\
$h_y,\deg_y$ & Global values on $\mathbb P^n(L)$ and $L^\times$; invariant under (5.3.3) & Verified linear representations\\
$\mathrm{vol}_y,\mu_y$ & Concrete measure objects and dependence require identification & Membership and nonmembership unclassified\\
\bottomrule
\end{tabularx}
\end{table}

%=====================================================================
\section{The Rigidity Theorem}\label{sec:rigidity}

\begin{theorem}[Rigidity of product-formula weights]\label{thm:rigidity}
Let $L$ be a number field and $(\gamma_v)_{v \in V_L}$ a family of real
exponents such that
\[
\prod_{v \in V_L} b_v(x)^{\gamma_v} = 1
\qquad \text{for all } x \in L^\times .
\]
Then $\gamma_v$ is constant in $v$. (We write $\gamma$ for the generic
exponent family of the theorem, reserving $\lambda$ for the source's
normalization exponents of Section~\ref{sec:types}.)
\end{theorem}

\begin{proof}
Set $\mu_v := \gamma_v - 1$. Subtracting the logarithm of the standard
product formula from the logarithm of the hypothesis gives
\[
\sum_{v \in V_L} \mu_v \log b_v(x) = 0
\qquad \text{for all } x \in L^\times,
\]
where for each fixed $x$ all but finitely many terms vanish.

\emph{(i) Archimedean component.} For units $u \in \OL^\times$ all finite
terms vanish, so $\sum_{\sigma\,\mathrm{arch}} \mu_\sigma \log b_\sigma(u) = 0$
for all $u$. By Dirichlet's unit theorem \cite[Ch.~I, \S7]{neukirch},
$\mathrm{Log}(\OL^\times)$ is a full lattice in the trace-zero hyperplane
$H = \{ y \in \mathbb{R}^{r_1 + r_2} : \sum_\sigma y_\sigma = 0 \}$: it lies
in $H$ because $\sum_\sigma \log b_\sigma(u) = \log|N(u)| = 0$, and its rank
is $r_1 + r_2 - 1 = \dim H$. The vector $(\mu_\sigma)_\sigma$ annihilates a
full lattice of $H$, hence all of $H$, and therefore lies in
$H^\perp = \mathbb{R}\cdot(1,\dots,1)$:
\[
\mu_\sigma = c \quad \text{for all archimedean } \sigma,
\text{ for some } c \in \mathbb{R}.
\]

\emph{(ii) Finite component.} Let $\mathfrak{p}$ be any prime ideal and $h$
the order of its class in the finite class group
\cite[Ch.~I, \S6]{neukirch}, so $\mathfrak{p}^h = (x_{\mathfrak{p}})$ is
principal. Then $\log b_{\mathfrak{q}}(x_{\mathfrak{p}}) = 0$ for all finite
$\mathfrak{q} \ne \mathfrak{p}$, while
$\log b_{\mathfrak{p}}(x_{\mathfrak{p}}) = -h \log N\mathfrak{p}$ and
\[
\sum_{\sigma} \log b_\sigma(x_{\mathfrak{p}})
= \log|N(x_{\mathfrak{p}})| = \log\bigl(N\mathfrak{p}^{\,h}\bigr)
= h \log N\mathfrak{p} .
\]
Substituting into the annihilation identity and using (i):
\[
\mu_{\mathfrak{p}} \cdot (-h \log N\mathfrak{p})
+ c \cdot (h \log N\mathfrak{p}) = 0
\quad\Longrightarrow\quad
\mu_{\mathfrak{p}} = c .
\]
As $\mathfrak{p}$ was arbitrary, $\mu_v = c$ for all $v$, i.e.,\
$\gamma_v = 1 + c$ is constant. Conversely, every constant family satisfies
the hypothesis.
\end{proof}

\begin{remark}[Status]\label{rem:status}
The theorem is elementary and lies in the orbit of the Artin--Whaples
axiomatics \cite{artinwhaples}; we state and prove it in full because the
series uses it as a typed test instrument, and because the consequence drawn
in Sections~\ref{sec:application}--\ref{sec:consequences} depends on its
exact hypotheses. No novelty claim is made for the theorem itself.
\end{remark}

\begin{corollary}[Unequal final place weights are excluded]\label{cor:noselect}
Under Theorem~\ref{thm:rigidity}'s same-field, full-place hypothesis,
only constant final exponent families satisfy the product formula.
In particular, prescribing $\gamma_v=j^2$ on a nonempty set
$S\subseteq V_L$ forces $\gamma_v=j^2$ at every place.
The constant global extension is admissible; extending the prescribed
values by unequal complementary weights is impossible.
\end{corollary}

\begin{lemma}[Finite-fiber redistribution]\label{lem:redistribution}
Let $\tau : V_L \to V_L$ have finite fibers (in particular, let $\tau$ be a
permutation), let $a_v > 0$, and suppose the family
$\bigl(b_{\tau(v)}^{\,a_v \gamma_v}\bigr)_{v}$ satisfies the product formula
for all $x \in L^\times$. Then the effective weights
$\beta_u := \sum_{v :\, \tau(v) = u} a_v \gamma_v$ are constant in $u$.
\end{lemma}

\begin{proof}
Reindexing the (for each fixed $x$ finite) sum,
$\sum_v a_v\gamma_v \log b_{\tau(v)}(x) = \sum_u \beta_u \log b_u(x) = 0$ for
all $x \in L^\times$ --- the hypothesis of Theorem~\ref{thm:rigidity} in the
effective weights. Finiteness of the fibers makes each $\beta_u$ a finite sum
and the reindexing legitimate.
\end{proof}

For infinite correspondences one would additionally need absolute
summability per fiber and a justified rearrangement; no such generality is
used in this paper, and none is needed for the audited passages, which move
no places.

%=====================================================================
\section{The Effective-Weight Quantity Class and the Exhaustion
Proposition}\label{sec:application}

\begin{definition}[Effective-weight quantity class]\label{def:geff}
Fix a domain $Z$. A family $Q=(Q_y:Z\to\mathbb R)_y$ lies in $\Geff$
if there are $y$-independent local functionals $q_w:Z\to\mathbb R$,
$w\in V_L$, such that
\[
Q_y(z)=\sum_{w\in V_L}\beta_{y,w}q_w(z),
\qquad Q_{\mathrm{ref}}(z):=\sum_{w\in V_L}q_w(z).
\]
Both sums must be finite or absolutely convergent, the first for every
$(y,z)$ and the reference sum for every $z$. The reference-sum requirement
also covers parameters with $\beta_y=0$.
\end{definition}

Membership requires this \emph{linear} representation; mere dependence on
$\beta_y$ is insufficient. For example, on a one-point domain take
$\beta_{y,w}=c_y$ at every place, with $c_y$ taking values $1$ and $2$,
and $Q_y=c_y^2$. A representation would give $Q_y=c_y S$, where
$S=\sum_wq_w$ is independent of $y$, forcing both $S=1$ and $S=2$.
This counterexample to the earlier membership gloss is part of the R18
correction, not a new claim about a source height.
The generic reference $Q_{\mathrm{ref}}$ need not be a classical
arithmetic quantity; that name requires a separately identified height or
degree and its domain.

\begin{lemma}[Factorization]\label{lem:factorization}
Assume the fixed-field, full-place weighted product formula
\[
\prod_v\rho_{y,v}(x)^{\lambda_{y,v}}=1
\qquad\text{for every }x\in L^\times.
\]
Then $\beta_{y,w}=c_y$ for all $w$, and every $Q\in\Geff$ satisfies
\[
Q_y=c_yQ_{\mathrm{ref}}.
\]
The scalar may depend on $y$. If (5.3.3) holds pointwise, then
$\beta_y=1$ and $Q_y=Q_{\mathrm{ref}}$.
Abstract real weights can have $c_y=0$; for the source's positive
$a_{y,v}$ and positive normalization exponents, $c_y>0$.
\end{lemma}

\begin{proof}
For every $x\in L^\times$, the hypothesis becomes
\[
0=\sum_v\lambda_{y,v}\log\rho_{y,v}(x)
=\sum_w\beta_{y,w}\log b_w(x).
\]
Theorem~\ref{thm:rigidity} yields $\beta_{y,w}=c_y$.
Substitution into Definition~\ref{def:geff} gives the result; the specified
convergence ensures the reference sum exists even when $c_y=0$.
\end{proof}

\begin{proposition}[Exhaustion on the stated arithmetic domains]\label{prop:exhaustion}
For the projective height of ATS~II($\tfrac12$), Definition~8.2.1
\cite{joshiIIhalf}, on $\mathbb P^n(L)$, and the weighted degree on
$L^\times$, use the fixed full standard family and the normalization
identity (5.3.3). Then
\begin{align*}
h_y([x_0:\dots:x_n])
&=\sum_v\lambda_{y,v}\log\max_i\rho_{y,v}(x_i)\\
&=h_{\mathrm{cl}}([x_0:\dots:x_n]),\\
\deg_y(x)&=\sum_v\lambda_{y,v}\log\rho_{y,v}(x)
=\deg_{\mathrm{cl}}(x).
\end{align*}
The degree equality has domain $x\in L^\times$.
\end{proposition}

\begin{proof}
On a nonzero representative tuple $(x_0,\dots,x_n)\in L^{n+1}\setminus\{0\}$,
positivity of $a_{y,v}$ and $\beta_{y,v}=1$ give
\begin{align*}
\sum_v\lambda_{y,v}\log\max_i\rho_{y,v}(x_i)
&=\sum_v\lambda_{y,v}a_{y,v}\log\max_i b_v(x_i)\\
&=\sum_v\log\max_i b_v(x_i).
\end{align*}
For each tuple, these sums have finite support. Changing the representative
by $u\in L^\times$ changes the standard sum by
$\sum_v\log b_v(u)=0$; the weighted sum has the same property by (5.3.4)
and (5.3.3), as reflected in Lemma~8.2.3 of the source.
Thus the equality descends to $\mathbb P^n(L)$.
For the degree the calculation omits the maximum and is valid on
$L^\times$.
\end{proof}

The local height functionals
$q_v(x_0,\dots,x_n)=\log\max_i b_v(x_i)$ are functions of representative
tuples, not individually well-defined projective functions.
They give a linear representation on the tuple domain.
For a representation on $\mathbb P^n(L)$, one may first fix a
$y$-independent representative for each point; the \emph{summed} height
is independent of that choice under the product formula.
The degree uses $q_v(x)=\log b_v(x)$ on $L^\times$.

\begin{remark}[Boundary and stabilization]\label{rem:boundary}
This is a statement about these identified domains and the linear class,
not every height or measure object in the ATS series.
In ATS~II($\tfrac12$) \S\S8.6--8.7, the stabilized height is evaluated on
$x\in L^\times$. Each summand $h_{\alpha\cdot y}(x)$ equals the same
classical value under the normalization identity, so its supremum equals
that value. No additional exchange of supremum and sum is needed for this
constant family. A general supremum of unclassified quantities has no
such conclusion. Log-volumes, hulls, Haar normalizations, and Jacobians
have neither membership nor nonmembership proved here; their concrete
representation is the open measure-pinning obligation R-E2.
\end{remark}

%=====================================================================
\section{Source-Bound Application to the Height Subchain}\label{sec:sourcebound}

Table~\ref{tab:substitution_audit} separates the direct substitution result
from source objects requiring an additional bridge. Page anchors refer to
the exact PDFs cited, not a claim to audit their entire contents.

\begin{table}[htbp]
\centering\small
\caption{Substitution results and unproved domain or object bridges.}\label{tab:substitution_audit}
\begin{tabularx}{\textwidth}{@{}p{0.26\textwidth}XX@{}}
\toprule
Source site & Consequence within the stated contract & Limit\\
\midrule
II($\tfrac12$) (5.3.3)--(5.3.4), p.~32 & $\beta_y=1$; the weighted product formula reduces to the standard formula & Same $L$ and full $V_L$\\
\addlinespace
II($\tfrac12$) Def.~8.2.1 / Lem.~8.2.3, p.~51 & $h_y=h_{\mathrm{cl}}$ on $\mathbb P^n(L)$ & Representative tuples must descend\\
\addlinespace
II($\tfrac12$) Rem.~8.2.2 / Prop.~8.3.1, p.~51 & The source extends to $\mathbb P^n(\overline L)$ and asserts nontrivial dependence there & Extension, valuation and base-change bridge unproved here; not refuted by the $L$-domain identity alone\\
\addlinespace
II($\tfrac12$) \S\S8.6--8.7, pp.~54--55 & On $L^\times$, every normalized summand equals $h_{\mathrm{cl}}$, hence so does the supremum & No conclusion for other stabilized objects\\
\addlinespace
III Rem.~8.9.2(2), p.~90 & Comparison of heights of points in $X(\overline L)$ using an ample line bundle is a separate source object & Bridge to the audited projective height remains open\\
\addlinespace
IV \S1.5, p.~17 & The term-by-term restriction concerns lower and upper log-volume comparisons across $y_0$ and $\varphi(y_0)$ & Not a universal prohibition of all local comparisons; no $\Geff$ classification\\
\addlinespace
II($\tfrac12$) Thm.~5.10.1(7) & Coordinate kernels can differ for nonproportional coefficients while their $L^\times$ pullbacks coincide & Conditional interpretation, not verification of the whole source theorem\\
\bottomrule
\end{tabularx}
\end{table}

\emph{Applicability gate H-PLACES.} The theorem uses one fixed number field,
all its archimedean and finite places, and all nonzero elements of that
field as tests. A source notation that ranges over nonarchimedean places
(such as ATS~III \S3.2(5)) cannot automatically be identified with this
$V_L$. Likewise, places of an extension $L'$ tested only on $L^\times$
require an explicit field-descent and place bridge. Neither indexing
identification is supplied by a coordinate analogy.

The source itself supplies the pointwise normalization used in the
$L$-domain substitution. This establishes the displayed height and degree
identities, but does not identify every operative height, line-bundle
quantity, or log-volume elsewhere in the series.

%=====================================================================
\section{Conditional Consequences (Three Typed Cases)}\label{sec:consequences}

The cases concern a verified quantity representation and an explicitly
matched domain, field, place family, and test set.
\begin{description}
\item[Case A (pointwise identity (5.3.3)).]\mbox{}\\
Then $\beta_y=1$ and every proved member of $\Geff$ equals its reference.
The identified height on $\mathbb P^n(L)$ and degree on $L^\times$
are classical and $y$-independent. Application to another operative
quantity requires identifying that object and its domain bridge.
\item[Case B (full weighted product formula only).]\mbox{}\\
Then $\beta_y=c_y\mathbf1$ and $Q_y=c_yQ_{\mathrm{ref}}$ for
$Q\in\Geff$. The scalar can vary with $y$, but final weights cannot vary
from place to place. For actual standard heights and degrees, this is
a global rescaling on their respective domains.
\item[Case C (membership unclassified).]\mbox{}\\
For the concrete log-volumes, hulls, and measures no linear
representation or obstruction to one is proved here.
The theorem still constrains effective weights whenever its weighted
product-formula hypothesis holds; it does not thereby classify these
quantity families. R-E2 remains open.
\end{description}

\begin{table}[htbp]
\centering\small
\caption{Typed consequence matrix.}\label{tab:case_verdicts}
\begin{tabularx}{\textwidth}{@{}p{0.16\textwidth}XXX@{}}
\toprule
Case & Effective weights & Global quantities & Derivation burden\\
\midrule
A: pointwise identity & $\beta_y=1$ & $Q_y=Q_{\mathrm{ref}}$ for proved members & Identify operative object and domain bridge\\
\addlinespace
B: full weighted formula & $\beta_y=c_y\mathbf1$ & $Q_y=c_yQ_{\mathrm{ref}}$ for proved members & Check weighted-formula hypothesis, indexing and factorization if another variation is intended\\
\addlinespace
C: unclassified quantity & Constrained if the same theorem hypotheses hold & No classification of the quantity from weights alone & Settle the concrete representation and measure-pinning bridge\\
\bottomrule
\end{tabularx}
\end{table}

\emph{Conditional consequence.} A place-selective metric variation cannot
be carried by the verified linear $L$-domain families under the full
weighted formula; the pointwise identity removes even their global
rescaling. Prop.~8.3.1 on $\mathbb P^n(\overline L)$, the ATS~III
line-bundle comparison, and ATS~IV's log-volume apparatus are separate
targets. Their operative object and transport bridges must be checked
before applying this conclusion. The standard product formula for the
fixed $b$ is not something these cases call on a source to disprove.
No global collapse, closure of the repair program, or exclusion of other
routes follows.

\emph{Falsification and revision criteria.}
\begin{enumerate}
\item A nonconstant effective-weight family satisfying exactly the same-field,
full-place product formula for every $x\in L^\times$ would contradict
Theorem~\ref{thm:rigidity}. A counterexample to factorization must also
satisfy the linear representation and convergence contract.
\item Evidence that an operative source height has another domain or
definition defeats an unproved identification step, not the theorem.
This is why the extension and line-bundle bridges remain separate.
\item A future source revision may supply new objects or bridges.
It would require a new audit of that version; it neither retroactively
proves an absent bridge nor refutes a statement tied to the cited version.
\end{enumerate}

%=====================================================================
\section[External Program Status (Not Part of Present Proofs)]{External Program Status\texorpdfstring{\\}{ }(Not Part of the Present Proofs)}\label{sec:notices}

Part~5 is published \cite{part5}; its publication does not discharge the
remaining mathematical bridges. R6-1 / REP-6-1 is a conditional algebraic
box-transport reconstruction: actual Haar normalization, tensor/index
compatibility and the operative volume comparison remain separate tasks.
A defect in an intermediate weighted model is not by itself a
counterexample to a Joshi theta-volume theorem.
R6-2 / REP-6-2 assumes V1 and an explicit $L'\to L$ place bridge.
R6-3 / REP-6-3 concerns uniqueness of the same-field normalized data,
not equality of all global coordinate hyperplanes.
These notices refer to the August program and its R18 corrections;
none is used in the proofs above.

The GNC source audit is cited in version 1.2 \cite{gnc}, and the
Part~4 ledger audit in version 0.3 \cite{part4}.
The certificate-form criterion is a separate instrument \cite{part2}.
Part~6's R-E2 measure-pinning obligation is not GNC-E2.
Likewise REP-6-1--3 are repair positions, not ledger F-R identifiers;
publication version 0.3 is not a claim to a new v1-F ledger acceptance.
The P5/Shadow Tomograph and full-chain obligations remain open.

%=====================================================================
\section{Claim Boundaries}\label{sec:claims}

Theorem~\ref{thm:rigidity} and Lemma~\ref{lem:redistribution} concern the
explicit fixed-field, full-place contract.
Lemma~\ref{lem:factorization} additionally requires a proved linear
representation; Proposition~\ref{prop:exhaustion} concerns
$\mathbb P^n(L)$ and $L^\times$.
Only these verified quantity/domain pairs support the negative statement
about residual place-selective variation.
The $\overline L$ extension, ATS~III line-bundle heights and ATS~IV
log-volumes have not been identified with that contract here.
No exhaustive absence of another metric carrier, proof of Case~C
nonmembership, closure of the repair program, or whole-source error
verdict is asserted. The coordinate-kernel interpretation is conditional
and does not certify the full nonconstancy assertion of source
Theorem~5.10.1(7). There is no proof or refutation of abc or IUT and no
whole-series mathematical acceptance.

The R18 inventory is the source of these corrections.
Independent manuscript and source checks are review evidence, not a
machine-checked Lean proof, a numerical experiment, or a new formal
verification result.

\section*{AI disclosure}
AI systems assisted in drafting, translation, and review of this working
paper. The corrective version uses the R18 inventory and independently
checked manuscript and selected source passages. The author remains
responsible for the text and its claims. These checks do not constitute
formal proof verification or mathematical acceptance of the whole ATS
series. No machine-checked proof or new numerical experiment is reported.


%=====================================================================
\begin{thebibliography}{99}
\raggedright

\bibitem{artinwhaples}
E.~Artin and G.~Whaples,
\emph{Axiomatic characterization of fields by the product formula for
valuations},
Bull.\ Amer.\ Math.\ Soc.\ \textbf{51} (1945), no.~7, 469--492.
DOI: \href{https://doi.org/10.1090/S0002-9904-1945-08383-9}{10.1090/S0002-9904-1945-08383-9}

\bibitem{neukirch}
J.~Neukirch,
\emph{Algebraic Number Theory},
Grundlehren der mathematischen Wissenschaften \textbf{322},
Springer, Berlin--Heidelberg, 1999.
ISBN 3-540-65399-6. DOI: \href{https://doi.org/10.1007/978-3-662-03983-0}{10.1007/978-3-662-03983-0}

\bibitem{joshiI}
K.~Joshi,
\emph{Construction of Arithmetic Teichmüller Spaces I},
arXiv:2106.11452v4 [math.AG], 2021 (revised 24 February 2025).

\bibitem{joshiII}
K.~Joshi,
\emph{Construction of Arithmetic Teichmüller Spaces II: Proof of a local
prototype of Mochizuki's Corollary 3.12},
arXiv:2303.01662v3 [math.AG], 2023 (revised 24 February 2025).

\bibitem{joshiIIhalf}
K.~Joshi,
\emph{Construction of Arithmetic Teichmüller Spaces II$\frac{1}{2}$:
Deformations of Number Fields},
arXiv:2305.10398v12 [math.AG], 2023 (revised 24 February 2025).

\bibitem{joshiIII}
K.~Joshi,
\emph{Construction of Arithmetic Teichmüller Spaces III: A `Rosetta Stone'
and a proof of Mochizuki's Corollary 3.12},
arXiv:2401.13508v4 [math.AG], 2024 (revised 24 February 2025).

\bibitem{joshiIV}
K.~Joshi,
\emph{Construction of Arithmetic Teichmüller Spaces IV: Proof of the
abc-conjecture},
arXiv:2403.10430v2 [math.AG], 2024 (revised 24 February 2025).

\bibitem{mochizukiIUT}
S.~Mochizuki,
\emph{Inter-universal Teichmüller Theory I--IV},
Publ.\ Res.\ Inst.\ Math.\ Sci.\ \textbf{57} (2021), no.~1/2, 3--723.

\bibitem{part1}
L.~Geiger,
\emph{The Comparison in IUTchIII, Corollary 3.12: From Attempted Repair to Structural Diagnosis},
Zenodo, 2026.
Concept DOI: \href{https://doi.org/10.5281/zenodo.19960781}{10.5281/zenodo.19960781}.
Cited version v6.4, 2026-10-01.
DOI: \href{https://doi.org/10.5281/zenodo.23081604}{10.5281/zenodo.23081604}.

\bibitem{part3}
L.~Geiger,
\emph{The Bridge Certificate: A Line-Level Audit Instrument for Inter-Theoretic Transfer Arguments, with Four Applications to the abc Literature},
Zenodo, 2026.
Concept DOI: \href{https://doi.org/10.5281/zenodo.21925803}{10.5281/zenodo.21925803}.
Cited version 0.3, 2026-10-01.
DOI: \href{https://doi.org/10.5281/zenodo.23090371}{10.5281/zenodo.23090371}.

\bibitem{part4}
L.~Geiger,
\emph{The Bridge Certificate Applied to the Joshi ATS Series: A Ledger Audit (DRAFT)},
Zenodo, 2026.
Concept DOI: \href{https://doi.org/10.5281/zenodo.21951155}{10.5281/zenodo.21951155}.
Cited version 0.3, 2026-10-02.
DOI: \href{https://doi.org/10.5281/zenodo.23092802}{10.5281/zenodo.23092802}.

\bibitem{gnc}
L.~Geiger,
\emph{Global Normalization in Joshi's Arithmetic Teichmüller Theory: What the Identification Carries --- and What Carries the Theorem (DRAFT)},
Zenodo, 2026.
Concept DOI: \href{https://doi.org/10.5281/zenodo.21924253}{10.5281/zenodo.21924253}.
Cited version 1.2, 2026-10-02.
DOI: \href{https://doi.org/10.5281/zenodo.23092166}{10.5281/zenodo.23092166}.

\bibitem{landscapeA}
L.~Geiger,
\emph{The abc Landscape: Obstructions, Failed Routes, and HCT Diagnostics (DRAFT)},
Zenodo, 2026.
Concept DOI: \href{https://doi.org/10.5281/zenodo.21916900}{10.5281/zenodo.21916900}.
Cited version 1.1, 2026-08-13.
DOI: \href{https://doi.org/10.5281/zenodo.21924338}{10.5281/zenodo.21924338}.

\bibitem{part2}
L.~Geiger,
\emph{The Bridge-Contract Criterion: A Form Audit for Transport Arguments, with an Instantiation in the abc Literature},
Zenodo, 2026.
Concept DOI: \href{https://doi.org/10.5281/zenodo.21960476}{10.5281/zenodo.21960476}.
Cited version 0.2, 2026-08-16.
DOI: \href{https://doi.org/10.5281/zenodo.21960477}{10.5281/zenodo.21960477}.

\bibitem{part5}
L.~Geiger,
\emph{The Joshi ATS Series under the Bridge-Certificate Standard: A Targeted Substance Audit and an All-Place Rigidity Classification of Effective Normalization Weights},
Zenodo, 2026.
Concept DOI: \href{https://doi.org/10.5281/zenodo.21956398}{10.5281/zenodo.21956398}.
Cited version 0.2, 2026-08-15.
DOI: \href{https://doi.org/10.5281/zenodo.21956399}{10.5281/zenodo.21956399}.

\end{thebibliography}

\end{document}
